\documentclass[12pt,a4paper]{article}

% ---- Paquetes ----
\usepackage[utf8]{inputenc}
\usepackage[T1]{fontenc}
\usepackage[spanish]{babel}
\usepackage{geometry}
\usepackage{hyperref}
\usepackage{xcolor}
\usepackage{listings}
\usepackage{enumitem}
\usepackage{float}
\usepackage{graphicx}
\usepackage{amsmath}
\usepackage{array}
\usepackage{longtable}
\usepackage{booktabs}
\usepackage{caption}

\geometry{margin=2.5cm}

% ---- Configuración de listings ----
\lstset{
    language=Python,
    basicstyle=\small\ttfamily,
    keywordstyle=\color{blue}\bfseries,
    stringstyle=\color{red},
    commentstyle=\color{green!50!black},
    frame=single,
    numbers=left,
    numberstyle=\tiny,
    breaklines=true,
    showstringspaces=false,
    tabsize=4,
    captionpos=b,
    backgroundcolor=\color{gray!5},
    literate=
        {á}{{\'a}}1 {é}{{\'e}}1 {í}{{\'i}}1 {ó}{{\'o}}1 {ú}{{\'u}}1
        {ñ}{{\~n}}1 {ü}{{\"u}}1 {Á}{{\'A}}1 {É}{{\'E}}1 {Í}{{\'I}}1
        {Ó}{{\'O}}1 {Ú}{{\'U}}1
}

\definecolor{codegreen}{rgb}{0,0.6,0}
\definecolor{codegray}{rgb}{0.5,0.5,0.5}
\definecolor{codepurple}{rgb}{0.58,0,0.82}
\definecolor{backcolour}{rgb}{0.95,0.95,0.92}

% ---- Metadata ----
\title{\textbf{Esquema de Validación Pydantic}\\para el Catálogo de Obras de Arte\\\large Proyecto Atrium -- Museo de Arte Contemporáneo}
\author{Documentación Técnica}
\date{\today}

\begin{document}

\maketitle
\thispagestyle{empty}
\newpage

\tableofcontents
\newpage

% ===================================================================
\section{Introducción}
% ===================================================================

El presente documento describe el diseño e implementación del esquema de validación
Pydantic para las obras de arte del sistema \textbf{Atrium}. El objetivo es
proporcionar una capa de validación robusta que maneje los distintos tipos de obras
(pintura, escultura, fotografía, cerámica y orfebrería) mediante un enfoque
\textbf{polimórfico} utilizando \texttt{Union} discriminada de Pydantic v2.

El archivo de implementación se encuentra en:
\begin{center}
\verb|app/schemas/artwork.py|
\end{center}

\subsection{Tecnologías utilizadas}

\begin{itemize}
    \item \textbf{Python 3.12+} con tipado estático (\texttt{typing}).
    \item \textbf{Pydantic v2} (\texttt{pydantic==2.13.4}) para validación y serialización.
    \item \textbf{FastAPI} (\texttt{fastapi==0.136.1}) para la capa web.
    \item \textbf{Beanie} (\texttt{beanie==2.1.0}) como ODM para MongoDB (usa Pydantic internamente).
\end{itemize}

% ===================================================================
\section{Estructura general del archivo}
% ===================================================================

El archivo \texttt{artwork.py} se organiza en seis secciones:

\begin{enumerate}
    \item Enumeraciones (\texttt{Genre}, \texttt{StatusObra}).
    \item Clase base (\texttt{ArtworkBase}) con atributos y validaciones comunes.
    \item Modelos específicos por género (\texttt{PinturaCreate}, \texttt{EsculturaCreate}, etc.).
    \item Unión discriminada para creación (\texttt{ArtworkCreate}).
    \item Modelos de respuesta con datos de base de datos (\texttt{PinturaResponse}, etc.).
    \item Unión discriminada para respuesta (\texttt{ArtworkResponse}).
    \item Filtros de búsqueda (\texttt{ArtworkFilter}).
\end{enumerate}

% ===================================================================
\section{Enumeraciones}
% ===================================================================

Las enumeraciones definen los valores fijos que puede tomar el discriminador
y el estado de una obra.

\subsection{Genre}

Define los cinco tipos de obra soportados. Cada valor es un string que viaja
en el JSON y actúa como discriminador polimórfico.

\begin{lstlisting}
class Genre(str, Enum):
    PINTURA = "pintura"
    ESCULTURA = "escultura"
    FOTOGRAFIA = "fotografia"
    CERAMICA = "ceramica"
    ORFEBRERIA = "orfebreria"
\end{lstlisting}

\subsection{StatusObra}

Representa el ciclo de vida de una obra dentro del sistema.

\begin{lstlisting}
class StatusObra(str, Enum):
    DISPONIBLE = "disponible"
    RESERVADA = "reservada"
    VENDIDA = "vendida"
\end{lstlisting}

Diagrama de estados:

\begin{center}
\ttfamily
\fbox{\textbf{Disponible}}
$\xrightarrow{\text{compra}}$
\fbox{\textbf{Reservada}}
$\xrightarrow{\text{concretar venta}}$
\fbox{\textbf{Vendida}}\\[4pt]
\hspace*{6.5cm}
$\xleftarrow[\text{cancelar (24h)}]{}$
\end{center}

% ===================================================================
\section{Clase base: ArtworkBase}
% ===================================================================

\subsection{Atributos comunes}

Contiene los campos que toda obra comparte, independientemente de su género.

\begin{lstlisting}
class ArtworkBase(BaseModel):
    nombre: str = Field(..., min_length=1, max_length=200)
    artista_id: str
    precio_venta: float = Field(..., gt=0)
    fecha_creacion: date
    foto: str
    descripcion: Optional[str] = Field(None, max_length=2000)
\end{lstlisting}

\begin{longtable}{lll}
\toprule
\textbf{Campo} & \textbf{Tipo} & \textbf{Validación} \\
\midrule
\endfirsthead
\toprule
\textbf{Campo} & \textbf{Tipo} & \textbf{Validación} \\
\midrule
\endhead
\bottomrule
\endfoot
\texttt{nombre} & \texttt{str} & 1--200 caracteres, obligatorio \\
\texttt{artista\_id} & \texttt{str} & Referencia al artista, obligatorio \\
\texttt{precio\_venta} & \texttt{float} & $> 0$, redondeado a 2 decimales \\
\texttt{fecha\_creacion} & \texttt{date} & No puede ser futura \\
\texttt{foto} & \texttt{str} & URL de la imagen, obligatorio \\
\texttt{descripcion} & \texttt{Optional[str]} & Máx. 2000 caracteres, opcional \\
\end{longtable}

\subsection{Validadores}

Se definen dos validadores de clase mediante \texttt{@field\_validator}:

\begin{lstlisting}
@field_validator("precio_venta")
@classmethod
def redondear_precio(cls, v: float) -> float:
    return round(v, 2)

@field_validator("fecha_creacion")
@classmethod
def no_futura(cls, v: date) -> date:
    if v > date.today():
        raise ValueError("La fecha de creacion no puede ser futura")
    return v
\end{lstlisting}

\begin{itemize}
    \item \texttt{redondear\_precio}: garantiza precisión monetaria de 2 decimales.
    \item \texttt{no\_futura}: evita fechas de creación inconsistentes.
\end{itemize}

\subsection{Herencia}

Todas las clases \texttt{*Create} heredan de \texttt{ArtworkBase}, por lo que
heredan automáticamente sus campos y validadores. No es necesario redefinirlos.

% ===================================================================
\section{Modelos específicos por género}
% ===================================================================

Cada género define una subclase de \texttt{ArtworkBase} que:
\begin{itemize}
    \item Fija el campo \texttt{genero} con un \texttt{Literal} de su valor
          correspondiente.
    \item Agrega los atributos particulares de ese tipo de obra.
\end{itemize}

\subsection{PinturaCreate}

\begin{lstlisting}
class PinturaCreate(ArtworkBase):
    genero: Literal[Genre.PINTURA] = Genre.PINTURA
    tecnica: str
    soporte: str
    alto_cm: float = Field(..., gt=0)
    ancho_cm: float = Field(..., gt=0)
\end{lstlisting}

\begin{center}
\begin{tabular}{lll}
\toprule
\textbf{Campo} & \textbf{Tipo} & \textbf{Descripción} \\
\midrule
\texttt{tecnica} & \texttt{str} & Óleo, acrílico, acuarela, témpera, etc. \\
\texttt{soporte} & \texttt{str} & Lienzo, madera, papel, cartón, etc. \\
\texttt{alto\_cm} & \texttt{float} & Altura en centímetros ($>0$) \\
\texttt{ancho\_cm} & \texttt{float} & Anchura en centímetros ($>0$) \\
\bottomrule
\end{tabular}
\end{center}

\subsection{EsculturaCreate}

\begin{lstlisting}
class EsculturaCreate(ArtworkBase):
    genero: Literal[Genre.ESCULTURA] = Genre.ESCULTURA
    material: str
    peso_kg: float = Field(..., gt=0)
    largo_cm: float = Field(..., gt=0)
    ancho_cm: float = Field(..., gt=0)
    profundidad_cm: float = Field(..., gt=0)
\end{lstlisting}

\begin{center}
\begin{tabular}{lll}
\toprule
\textbf{Campo} & \textbf{Tipo} & \textbf{Descripción} \\
\midrule
\texttt{material} & \texttt{str} & Bronce, mármol, madera, hierro, etc. \\
\texttt{peso\_kg} & \texttt{float} & Peso en kilogramos ($>0$) \\
\texttt{largo\_cm} & \texttt{float} & Largo en centímetros ($>0$) \\
\texttt{ancho\_cm} & \texttt{float} & Ancho en centímetros ($>0$) \\
\texttt{profundidad\_cm} & \texttt{float} & Profundidad en centímetros ($>0$) \\
\bottomrule
\end{tabular}
\end{center}

\subsection{FotografiaCreate}

\begin{lstlisting}
class FotografiaCreate(ArtworkBase):
    genero: Literal[Genre.FOTOGRAFIA] = Genre.FOTOGRAFIA
    tecnica: str
    papel: Optional[str] = None
    alto_cm: float = Field(..., gt=0)
    ancho_cm: float = Field(..., gt=0)
    edicion: Optional[int] = Field(None, ge=1)
\end{lstlisting}

\begin{center}
\begin{tabular}{lll}
\toprule
\textbf{Campo} & \textbf{Tipo} & \textbf{Descripción} \\
\midrule
\texttt{tecnica} & \texttt{str} & Digital, analógica \\
\texttt{papel} & \texttt{Optional[str]} & Tipo de soporte fotográfico \\
\texttt{alto\_cm} & \texttt{float} & Altura en centímetros ($>0$) \\
\texttt{ancho\_cm} & \texttt{float} & Anchura en centímetros ($>0$) \\
\texttt{edicion} & \texttt{Optional[int]} & Número de edición/tirada ($\ge 1$) \\
\bottomrule
\end{tabular}
\end{center}

\subsection{CeramicaCreate}

\begin{lstlisting}
class CeramicaCreate(ArtworkBase):
    genero: Literal[Genre.CERAMICA] = Genre.CERAMICA
    tipo_arcilla: str
    tecnica_coccion: str
    peso_kg: float = Field(..., gt=0)
    alto_cm: float = Field(..., gt=0)
    ancho_cm: float = Field(..., gt=0)
    profundidad_cm: float = Field(..., gt=0)
    esmaltado: Optional[bool] = None
\end{lstlisting}

\begin{center}
\begin{tabular}{lll}
\toprule
\textbf{Campo} & \textbf{Tipo} & \textbf{Descripción} \\
\midrule
\texttt{tipo\_arcilla} & \texttt{str} & Porcelana, gres, loza, etc. \\
\texttt{tecnica\_coccion} & \texttt{str} & Bizcocho, esmaltado, raku, etc. \\
\texttt{peso\_kg} & \texttt{float} & Peso en kilogramos ($>0$) \\
\texttt{alto\_cm} & \texttt{float} & Altura en centímetros ($>0$) \\
\texttt{ancho\_cm} & \texttt{float} & Anchura en centímetros ($>0$) \\
\texttt{profundidad\_cm} & \texttt{float} & Profundidad en centímetros ($>0$) \\
\texttt{esmaltado} & \texttt{Optional[bool]} & Indica si tiene esmaltado \\
\bottomrule
\end{tabular}
\end{center}

\subsection{OrfebreriaCreate}

\begin{lstlisting}
class OrfebreriaCreate(ArtworkBase):
    genero: Literal[Genre.ORFEBRERIA] = Genre.ORFEBRERIA
    material: str
    tecnica: str
    peso_g: float = Field(..., gt=0)
    alto_cm: float = Field(..., gt=0)
    ancho_cm: float = Field(..., gt=0)
    profundidad_cm: float = Field(..., gt=0)
    quilates: Optional[float] = Field(None, ge=0, le=24)
\end{lstlisting}

\begin{center}
\begin{tabular}{lll}
\toprule
\textbf{Campo} & \textbf{Tipo} & \textbf{Descripción} \\
\midrule
\texttt{material} & \texttt{str} & Oro, plata, bronce, cobre, etc. \\
\texttt{tecnica} & \texttt{str} & Filigrana, repujado, fundición, etc. \\
\texttt{peso\_g} & \texttt{float} & Peso en gramos ($>0$) \\
\texttt{alto\_cm} & \texttt{float} & Altura en centímetros ($>0$) \\
\texttt{ancho\_cm} & \texttt{float} & Anchura en centímetros ($>0$) \\
\texttt{profundidad\_cm} & \texttt{float} & Profundidad en centímetros ($>0$) \\
\texttt{quilates} & \texttt{Optional[float]} & Pureza del material (0--24) \\
\bottomrule
\end{tabular}
\end{center}

% ===================================================================
\section{Polimorfismo: ArtworkCreate}
% ===================================================================

\subsection{La unión discriminada}

El núcleo del diseño polimórfico es \texttt{ArtworkCreate}, un type alias
que combina los 5 modelos mediante una \texttt{Union} anotada con un
discriminador:

\begin{lstlisting}
ArtworkCreate = Annotated[
    Union[
        PinturaCreate,
        EsculturaCreate,
        FotografiaCreate,
        CeramicaCreate,
        OrfebreriaCreate,
    ],
    Field(discriminator="genero")
]
\end{lstlisting}

\subsection{Funcionamiento interno}

Cuando Pydantic recibe un JSON para validar contra \texttt{ArtworkCreate}:

\begin{enumerate}
    \item Lee el campo \texttt{genero} del JSON entrante.
    \item Busca qué modelo de la \texttt{Union} tiene un \texttt{Literal} con
          ese valor.
    \item Instancia y valida ese modelo específico.
    \item Si el valor de \texttt{genero} no coincide con ningún modelo,
          lanza un error de validación.
\end{enumerate}

\begin{center}
\begin{minipage}{0.9\textwidth}
\footnotesize
\begin{verbatim}
   JSON entrante
  {..."genero": "escultura", "material": "bronce"...}
         |
         v
  +------------------+
  |  genero = ?      |  <-- discriminador
  +------------------+
         |
         | "escultura"
         v
  +------------------+
  | EsculturaCreate  |  <-- Pydantic selecciona este modelo
  +------------------+      automaticamente
\end{verbatim}
\end{minipage}
\captionof{figure}{Proceso de discriminación automática en ArtworkCreate}
\end{center}

\subsection{Uso en FastAPI}

\begin{lstlisting}
@app.post("/obras")
async def crear_obra(obra: ArtworkCreate):
    # FastAPI + Pydantic resuelven el tipo concreto automaticamente
    if obra.genero == Genre.PINTURA:
        obra.tecnica   # disponible solo en PinturaCreate
    elif obra.genero == Genre.ESCULTURA:
        obra.material  # disponible solo en EsculturaCreate
\end{lstlisting}

Sin el discriminador, FastAPI intentaría validar el JSON contra todos los
modelos de la \texttt{Union} secuencialmente hasta encontrar uno que funcione
(\emph{try-each}), lo cual es ineficiente y puede producir errores confusos.
El discriminador resuelve el tipo en \textbf{O(1)}.

\subsection{Ventajas del enfoque}

\begin{itemize}
    \item \textbf{Type safety}: el tipado estático (IDE, mypy) sabe qué campos
          tiene cada variante.
    \item \textbf{Validación específica}: cada género valida solo sus campos
          relevantes.
    \item \textbf{Extensible}: agregar un nuevo género requiere crear una
          subclase y añadirla a la \texttt{Union}.
    \item \textbf{Documentación viva}: FastAPI genera automáticamente el
          esquema OpenAPI con todos los modelos y el discriminador.
\end{itemize}

% ===================================================================
\section{Modelos de respuesta}
% ===================================================================

\subsection{ArtworkResponseBase}

Contiene campos que solo existen en la base de datos, no en la petición:

\begin{lstlisting}
class ArtworkResponseBase(BaseModel):
    id: str = Field(..., alias="_id")
    status: StatusObra = Field(default=StatusObra.DISPONIBLE)
    created_at: datetime = Field(default_factory=datetime.utcnow)
    updated_at: Optional[datetime] = None

    model_config = ConfigDict(
        populate_by_name=True,
        from_attributes=True,
    )
\end{lstlisting}

\begin{longtable}{lll}
\toprule
\textbf{Campo} & \textbf{Tipo} & \textbf{Descripción} \\
\midrule
\endfirsthead
\toprule
\textbf{Campo} & \textbf{Tipo} & \textbf{Descripción} \\
\midrule
\endhead
\bottomrule
\endfoot
\texttt{id} & \texttt{str (alias \_id)} & ID de MongoDB \\
\texttt{status} & \texttt{StatusObra} & Estado actual (default: disponible) \\
\texttt{created\_at} & \texttt{datetime} & Fecha de creación del registro \\
\texttt{updated\_at} & \texttt{Optional[datetime]} & Última modificación \\
\end{longtable}

\subsection{Herencia múltiple para responses}

Cada modelo de respuesta combina el Create correspondiente con
\texttt{ArtworkResponseBase}:

\begin{lstlisting}
class PinturaResponse(PinturaCreate, ArtworkResponseBase):
    model_config = ConfigDict(
        populate_by_name=True,
        from_attributes=True,
    )
\end{lstlisting}

Esto significa que \texttt{PinturaResponse} tiene:
\begin{itemize}
    \item Todos los campos de \texttt{PinturaCreate} (nombre, técnica, etc.).
    \item Todos los campos de \texttt{ArtworkResponseBase} (id, status, etc.).
\end{itemize}

\subsection{ArtworkResponse}

Análogo a \texttt{ArtworkCreate}, pero con los modelos de respuesta:

\begin{lstlisting}
ArtworkResponse = Annotated[
    Union[
        PinturaResponse, EsculturaResponse, FotografiaResponse,
        CeramicaResponse, OrfebreriaResponse,
    ],
    Field(discriminator="genero")
]
\end{lstlisting}

\subsection{ConfigDict: from\_attributes}

\texttt{from\_attributes=True} permite instanciar el modelo desde un objeto
arbitrario (no solo un diccionario), lo cual es esencial para integrar con
Beanie/MongoDB:

\begin{lstlisting}
# Desde un documento Beanie recuperado de MongoDB:
doc = await Artwork.find_one(Artwork.id == obra_id)
return PinturaResponse.model_validate(doc)  # funciona!
\end{lstlisting}

\section{Comparativa Create vs Response}

\begin{center}
\begin{tabular}{l|c|c}
\toprule
\textbf{Característica} & \textbf{ArtworkCreate} & \textbf{ArtworkResponse} \\
\midrule
Incluye \texttt{id} (MongoDB) & No & Sí \\
Incluye \texttt{status} & No & Sí (default: disponible) \\
Incluye timestamps & No & Sí \\
Validación de entrada & Sí & No (solo salida) \\
\texttt{from\_attributes} & No & Sí \\
Uso típico & POST /obras & GET, PUT /obras \\
\bottomrule
\end{tabular}
\end{center}

% ===================================================================
\section{Filtros de búsqueda: ArtworkFilter}
% ===================================================================

Para los endpoints de consulta, se define \texttt{ArtworkFilter} con campos
opcionales que permiten construir dinámicamente la query de MongoDB:

\begin{lstlisting}
class ArtworkFilter(BaseModel):
    genero: Optional[Genre] = None
    artista_id: Optional[str] = None
    precio_min: Optional[float] = Field(None, ge=0)
    precio_max: Optional[float] = Field(None, ge=0)
    status: Optional[StatusObra] = None
\end{lstlisting}

\begin{center}
\begin{tabular}{lll}
\toprule
\textbf{Campo} & \textbf{Tipo} & \textbf{Uso en MongoDB} \\
\midrule
\texttt{genero} & \texttt{Optional[Genre]} & \texttt{\{genero: valor\}} \\
\texttt{artista\_id} & \texttt{Optional[str]} & \texttt{\{artista\_id: valor\}} \\
\texttt{precio\_min} & \texttt{Optional[float]} & \texttt{\{precio\_venta: \{\$gte: valor\}\}} \\
\texttt{precio\_max} & \texttt{Optional[float]} & \texttt{\{precio\_venta: \{\$lte: valor\}\}} \\
\texttt{status} & \texttt{Optional[StatusObra]} & \texttt{\{status: valor\}} \\
\bottomrule
\end{tabular}
\end{center}

Ejemplo de uso en un endpoint GET:

\begin{lstlisting}
@app.get("/obras")
async def listar_obras(
    filters: Annotated[ArtworkFilter, Query()]
):
    query = {}
    if filters.genero:
        query["genero"] = filters.genero.value
    if filters.precio_min or filters.precio_max:
        query["precio_venta"] = {}
        if filters.precio_min:
            query["precio_venta"]["$gte"] = filters.precio_min
        if filters.precio_max:
            query["precio_venta"]["$lte"] = filters.precio_max
    if filters.status:
        query["status"] = filters.status.value
    ...
\end{lstlisting}

% ===================================================================
\section{Casos de uso}
% ===================================================================

\subsection{Creación de una pintura}

\begin{lstlisting}
# Request POST /obras
{
    "nombre": "La noche estrellada",
    "artista_id": "abc123",
    "precio_venta": 15000.50,
    "fecha_creacion": "1889-06-01",
    "foto": "https://ejemplo.com/noche.jpg",
    "genero": "pintura",
    "tecnica": "oleo",
    "soporte": "lienzo",
    "alto_cm": 73.7,
    "ancho_cm": 92.1
}
\end{lstlisting}

Pydantic:
\begin{enumerate}
    \item Lee \texttt{genero: "pintura"}.
    \item Selecciona \texttt{PinturaCreate}.
    \item Valida todos los campos de \texttt{PinturaCreate}.
    \item Redondea \texttt{precio\_venta} a 2 decimales.
    \item Verifica que \texttt{fecha\_creacion} no sea futura.
\end{enumerate}

\subsection{Error de validación}

\begin{lstlisting}
# Request POST /obras (incompleto)
{
    "nombre": "El pensador",
    "genero": "escultura",
    "material": "bronce"
    # Falta: precio_venta, fecha_creacion, foto,
    #        peso_kg, largo_cm, ancho_cm, profundidad_cm
}
\end{lstlisting}

Respuesta (\texttt{422 Unprocessable Entity}):

\begin{lstlisting}
{
    "detail": [
        {"loc": ["body", "precio_venta"],
         "msg": "Field required", "type": "missing"},
        {"loc": ["body", "fecha_creacion"],
         "msg": "Field required", "type": "missing"},
        {"loc": ["body", "foto"],
         "msg": "Field required", "type": "missing"},
        {"loc": ["body", "peso_kg"],
         "msg": "Field required", "type": "missing"},
        {"loc": ["body", "largo_cm"],
         "msg": "Field required", "type": "missing"},
        {"loc": ["body", "ancho_cm"],
         "msg": "Field required", "type": "missing"},
        {"loc": ["body", "profundidad_cm"],
         "msg": "Field required", "type": "missing"}
    ]
}
\end{lstlisting}

\subsection{Búsqueda con filtros}

\begin{lstlisting}
# Request GET /obras?genero=escultura&precio_min=1000&precio_max=50000
\end{lstlisting}

FastAPI parsea los query parameters en \texttt{ArtworkFilter}:
\begin{lstlisting}
ArtworkFilter(genero=Genre.ESCULTURA,
              precio_min=1000.0,
              precio_max=50000.0)
\end{lstlisting}

Query MongoDB resultante:
\begin{lstlisting}
{
    "genero": "escultura",
    "precio_venta": {"$gte": 1000, "$lte": 50000}
}
\end{lstlisting}

% ===================================================================
\section{Diagrama completo del flujo}
% ===================================================================

\begin{center}
\begin{minipage}{0.9\textwidth}
\footnotesize
\begin{verbatim}
  +-------------------+
  |  Cliente (JSON)   |
  +-------------------+
           |
           v
  +-------------------+
  | FastAPI Endpoint   |
  +-------------------+
           |
           v
  +-------------------+
  | genero = ?        |  <- discriminador
  +-------------------+
     /              \
    /                \
   v                  v
+----------+     +----------+
| Valida-  |     | Error    |
| ción     |     | 422      |
| especí-  |     +----------+
| fica     |
+----------+
    |
    v
+----------+
| MongoDB  |
| (Beanie) |
+----------+
    |
    v
+----------+
| Respues- |
| ta dis-  |
| criminada|
+----------+
    |
    |  (respuesta JSON)
    v
  Cliente
\end{verbatim}
\end{minipage}
\captionof{figure}{Flujo completo de validación y persistencia}
\end{center}

% ===================================================================
\section{Extensibilidad}
% ===================================================================

Para agregar un nuevo género (ej. \texttt{video}), los pasos son:

\begin{enumerate}
    \item Agregar el valor a \texttt{Genre}:
\begin{lstlisting}
class Genre(str, Enum):
    ...
    VIDEO = "video"
\end{lstlisting}
    \item Crear el modelo específico:
\begin{lstlisting}
class VideoCreate(ArtworkBase):
    genero: Literal[Genre.VIDEO] = Genre.VIDEO
    duracion_seg: int = Field(..., gt=0)
    resolucion: str
    formato: str
\end{lstlisting}
    \item Crear el modelo de respuesta:
\begin{lstlisting}
class VideoResponse(VideoCreate, ArtworkResponseBase):
    model_config = ConfigDict(populate_by_name=True,
                              from_attributes=True)
\end{lstlisting}
    \item Agregarlo a las uniones:
\begin{lstlisting}
ArtworkCreate = Annotated[
    Union[..., VideoCreate], Field(discriminator="genero")]
ArtworkResponse = Annotated[
    Union[..., VideoResponse], Field(discriminator="genero")]
\end{lstlisting}
\end{enumerate}

Sin modificar ninguna otra parte del código.

% ===================================================================
\section{Conclusión}
% ===================================================================

El diseño implementado proporciona:

\begin{itemize}
    \item \textbf{Validación tipo-específica}: cada género valida exactamente
          los campos que le corresponden.
    \item \textbf{Polimorfismo eficiente}: el discriminador resuelve el tipo
          en O(1), sin prueba secuencial.
    \item \textbf{Segregación Create/Response}: los modelos de creación no
          exponen campos internos de la base de datos.
    \item \textbf{Documentación automática}: FastAPI genera el esquema
          OpenAPI completo a partir de las anotaciones.
    \item \textbf{Extensibilidad}: agregar un nuevo género requiere solo
          añadir clases, sin modificar la lógica existente.
\end{itemize}

\end{document}
